
\SGAStatement{\SGAInlineRefTarget{sga3:xvii:proposition:3-10}Proposition 3.10}{Let \(G\) be an algebraic group
defined over an algebraically closed field \(k\).  The following
properties are equivalent:
\begin{enumerate}
\item[(i)] \SGARefTarget{sga3:xvii:proposition:3:10:item:i}
  Every homomorphism from a group \(M\) of multiplicative type into
  \(G\) is trivial.
\item[(i bis)] \SGARefTarget{sga3:xvii:proposition:3:10:item:i:bis}
  \(G\) has no nontrivial subgroup of multiplicative type.
\item[(ii)] \SGARefTarget{sga3:xvii:proposition:3:10:item:ii}
  \begin{enumerate}
  \item[(a)] \SGARefTarget{sga3:xvii:proposition:3:10:item:ii:a}
    If \(p=0\), then \(G(k)\) has no point of finite order other
    than \(e\).
  \item[(b)] \SGARefTarget{sga3:xvii:proposition:3:10:item:ii:b}
    If \(p\ne0\), then:
    \begin{itemize}
    \item for every prime number \(q\ne p\), if \(x\in G(k)\) and
      \(x^q=e\), then \(x=e\);
    \item \SGARefTarget{sga3:xvii:proposition:3:10:item:ii:b:lie}
      for every \(X\in\mathfrak g=\operatorname{Lie}G\) such that
      \(X^{(p)}=X\), one has \(X=0\).
    \end{itemize}
  \end{enumerate}
\item[(ii bis)] \SGARefTarget{sga3:xvii:proposition:3:10:item:ii:bis}
  \begin{enumerate}
  \item[(a)] \SGARefTarget{sga3:xvii:proposition:3:10:item:ii:bis:a}
    If \(p=0\), the condition is the same as in \SGARefLink{sga3:xvii:proposition:3:10:item:ii}{(ii)}\SGARefLink{sga3:xvii:appendix:A}{(a)}.
  \item[(b)] \SGARefTarget{sga3:xvii:proposition:3:10:item:ii:bis:b}
    If \(p\ne0\), then:
    \begin{itemize}
    \item every finite subgroup of \(G(k)\) is a \(p\)-group;
    \item \SGARefTarget{sga3:xvii:proposition:3:10:item:ii:bis:b:lie}
      for every \(X\in\mathfrak g=\operatorname{Lie}G\) such that
      \(X^{(p)}=X\), one has \(X=0\).
    \end{itemize}
  \end{enumerate}
\end{enumerate}}
\SGARefTarget{sga3:xvii:proposition:3:10}

\SGAProof{Proof}
The equivalence \SGARefLink{sga3:xvii:proposition:3:10:item:i}{(i)}\(\Leftrightarrow\)(i bis) follows because the
image of a group of multiplicative type is of multiplicative type
(\SGARefLink{sga3:ix:proposition:2:7}{Exposé IX, 2.7}).

The equivalence \SGARefLink{sga3:ix:proposition:2:7:item:ii:01}{(ii)}\(\Leftrightarrow\)(ii bis) follows from the
following observation.  If an ordinary finite group \(H\) has order
that is not a power of \(p\), then there exist a prime number
\(q\ne p\) and an element \(x\in H\), distinct from \(e\), such
that \(x^q=e\) (Sylow's theorem; cf. J.-P. Serre,
\emph{Corps locaux}, Chapter IX, \S2).

The implication \SGARefLink{sga3:xvii:proposition:3:10:item:i}{(i)}\(\Rightarrow\)\SGARefLink{sga3:xvii:proposition:3:10:item:ii}{(ii)} follows from the next lemma.

\SGAStatement{\SGAInlineRefTarget{sga3:xvii:lemma:3-11}Lemma 3.11}{Let \(G\) be an algebraic group defined
over a field \(k\).
\begin{enumerate}
\item[(a)] \SGARefTarget{sga3:xvii:lemma:3:11:item:a}
  If \(k\) contains the \(n\)-th roots of unity, where \(n\) is
  prime to \(p\), then
  \[
    \operatorname{Hom}_{k\text{-}\mathrm{gr}}(\mu_n,G)
      \simeq
    \operatorname{Hom}_{k\text{-}\mathrm{gr}}
      (\mathbb Z/n\mathbb Z,G)
      \simeq {}_nG(k),
  \]
  where \({}_nG(k)\) denotes the points of order \(n\) in \(G(k)\).
\item[(b)] \SGARefTarget{sga3:xvii:lemma:3:11:item:b}
  If \(p>0\), then
  \[
    \operatorname{Hom}_{k\text{-}\mathrm{gr}}(\mu_p,G)
      \simeq
    \{X\in\mathfrak g=\operatorname{Lie}G\mid X^{(p)}=X\}.
  \]
\end{enumerate}}
\SGARefTarget{sga3:xvii:lemma:3:11}

Indeed, for \SGARefLink{sga3:xvii:appendix:A}{(a)}, note that \(\mu_n\) is then isomorphic over \(k\)
to \(\mathbb Z/n\mathbb Z\); assertion \SGARefLink{sga3:xvii:section:b}{(b)} follows from Appendix II,
\SGARefLink{sga3:xvii:proposition:2:1}{2.1}.

For \SGARefLink{sga3:xvii:proposition:2:1:item:ii}{(ii)}\(\Rightarrow\)(i bis), \SGARefLink{sga3:xvii:lemma:3-11}{Lemma 3.11} shows that \SGARefLink{sga3:xvii:proposition:2:1:item:ii}{(ii)} is
equivalent to the assertion that, for every prime number \(r\), the
group \(G\) contains no subgroup \(\mu_r\).  The implication now
follows from:

\SGAStatement{Lemma 3.12}{Let \(G\) be a diagonalizable
\(S\)-group, of finite type over \(S\), and distinct from the
trivial group.  There exist a prime number \(r\) and a subgroup of
\(G\) isomorphic to \((\mu_r)_S\).}
\SGARefTarget{sga3:xvii:lemma:3:12}

Write \(G=D_S(M)\), where \(M\) is a finitely generated abelian
group, hence an extension of a free group \(M''\) by a finite group
\(M'\).  If \(M''\ne0\), it is clear that \(M\) has quotients
isomorphic to \(\mathbb Z/r\mathbb Z\) for every integer \(r\).  If
\(M''=0\), then \(M'\) has a quotient isomorphic to
\(\mathbb Z/r\mathbb Z\) for every prime number \(r\) dividing the
order of \(M\).  The lemma follows by duality.

We saw in \SGARefLink{sga3:xvii:proposition:2:4}{Proposition 2.4} that a unipotent group satisfies the
equivalent conditions of \SGARefLink{sga3:xvii:proposition:3-10}{3.10}.  The purpose of the next section is
to prove the converse.

\SGASection{4}{A Characterization of Unipotent Groups}
\SGARefTarget{sga3:xvii:section:4}

As announced, we shall prove that an algebraic group \(G\) defined
over an algebraically closed field \(k\), and containing no
nontrivial subgroup of multiplicative type, is unipotent.  In fact,
it is enough that \(G\) contain none of certain particular
``elementary'' subgroups of multiplicative type, which depend on the
hypotheses imposed on \(G\).  Before stating the general theorem,
let us study some special cases in detail.
